% META: {"id":"1SPE-GEOREP-RE-C3","chapitre":"1SPE-GEOMETRIE-REPEREE","type_objet":"remediation","capacites_codes":["C3"],"status":"approved","origin":"generated","status_history":["generated","audited","accepted","approved"],"release_acceptance":"RELEASE_OWNER_BATCH_ACCEPTANCE","acceptance_closure_digest":"sha256:9b3ccf9a81c5520fbb7e03b7b2d3e7bf2057a02834b6d908bf3be5f49f3b553f","acceptance_receipt_digest":"sha256:4fad86f0282e0fa4e0419cca1ad6379ae7af5a280c04a4a90880f90a1c044242"}

%---------------------------------------------------------------
% FICHE DE REMÉDIATION — C3
% Équation du cercle
%---------------------------------------------------------------

\begin{exercice}{1SPE-GEOREP-RE-C3-EX1}{1}{8}
Écrire l'équation du cercle de centre $\Omega(1 ; -4)$ et de rayon $3$.
\end{exercice}

% BEGIN-VERIFY
% assert 3**2 == 9
% END-VERIFY

\begin{corrige}{1SPE-GEOREP-RE-C3-EX1}

$(x - 1)^2 + (y + 4)^2 = 9$.

\end{corrige}

\begin{exercice}{1SPE-GEOREP-RE-C3-EX2}{1}{10}
Identifier le centre et le rayon du cercle $x^2 + y^2 + 8x - 2y - 8 = 0$.
\end{exercice}

% BEGIN-VERIFY
% # A=8, B=-2, C=-8. Centre (-4, 1). r^2 = 16+1+8 = 25
% assert 16 + 1 + 8 == 25
% END-VERIFY

\coupDePouce{1}{Complète le carré : $x^2 + 8x = (x+4)^2 - 16$.}

\begin{corrige}{1SPE-GEOREP-RE-C3-EX2}

$x^2 + 8x = (x+4)^2 - 16$ et $y^2 - 2y = (y-1)^2 - 1$.

$(x+4)^2 + (y-1)^2 = 25$. Centre $(-4 ; 1)$, rayon $5$.

\end{corrige}

\begin{exercice}{1SPE-GEOREP-RE-C3-EX3}{1}{8}
Le point $P(5 ; 0)$ est-il sur le cercle $(x - 2)^2 + (y - 4)^2 = 25$ ?
\end{exercice}

% BEGIN-VERIFY
% assert (5-2)**2 + (0-4)**2 == 25
% END-VERIFY

\begin{corrige}{1SPE-GEOREP-RE-C3-EX3}

$(5 - 2)^2 + (0 - 4)^2 = 9 + 16 = 25$. Donc $P$ est sur le cercle.

\end{corrige}
